\ifdefined\gkver\else\def\gkver{student}\fi
\documentclass[\gkver]{gaokaozhenti}
\begin{document}

\chapter{2024年安徽}

\begin{enumerate}
\item
%% number: 1
%% paperName: 2024年安徽高考第 1 题 4 分
%% typeId: 1
%% type: 单选题
%% chapter: 原子和原子核
%% point: 玻尔模型
%% method: 无
%% score: 4
%% degree: 850
%% duplicateId: 0
%% body:
大连相干光源是我国第一台高增益自由电子激光用户装置，其激光辐射所应用的玻尔原子理论很好地解释了氢原子的光谱特征。图为氢原子的能级示意图，已知紫外光的光子能量大于 $3.11 \UeV$，当大量处于 $n = 3$ 能级的氢原子向低能级跃迁时，辐射不同频率的紫外光有\xzanswer{B}

\onepicture[w=3.20cm, align=right]{\includesvg{figs/01}}

\fourchoices[answer=B]
{1 种}
{2 种}
{3 种}
{4 种}

%% answer: B
%% memo:
\memoanswer{
大量处于 $n = 3$ 能级的氢原子向低能级跃迁时，能够辐射出不同频率的种类为 $C_{3}^{2} = 3$ 种。

辐射出光子的能量分别为

$\Delta E_{1} = E_{3} - E_{1} = -1.51 \UeV - (-13.6 \UeV) = 12.09 \UeV$

$\Delta E_{2} = E_{3} - E_{2} = -1.51 \UeV - (-3.4 \UeV) = 1.89 \UeV$

$\Delta E_{3} = E_{2} - E_{1} = -3.4 \UeV - (-13.6 \UeV) = 10.2 \UeV$

其中 $\Delta E_{1} > 3.11 \UeV$，$\Delta E_{2} < 3.11 \UeV$，$\Delta E_{3} > 3.11 \UeV$

所以辐射不同频率的紫外光有 2 种。

故选 B。
}

\item
%% number: 2
%% paperName: 2024年安徽高考第 2 题 4 分
%% typeId: 1
%% type: 单选题
%% chapter: 功和能
%% point: 动能定理
%% method: 无
%% score: 4
%% degree: 850
%% duplicateId: 0
%% body:
某同学参加户外拓展活动，遵照安全规范，坐在滑板上，从高为 $h$ 的粗糙斜坡顶端由静止下滑，至底端时速度为 $v$。已知人与滑板的总质量为 $m$，可视为质点．重力加速度大小为 $g$，不计空气阻力．则此过程中人与滑板克服摩擦力做的功为\xzanswer{D}

\fourchoices[answer=D]
{$mgh$}
{$\frac{1}{2}mv^{2}$}
{$mgh + \frac{1}{2}mv^{2}$}
{$mgh - \frac{1}{2}mv^{2}$}

%% answer: D
%% memo:
\memoanswer{
人在下滑的过程中，由动能定理可得

$mgh - W_{f} = \frac{1}{2}mv^{2} - 0$

可得此过程中人与滑板克服摩擦力做的功为 $W_{f} = mgh - \frac{1}{2}mv^{2}$

故选 D。
}

\item
%% number: 3
%% paperName: 2024年安徽高考第 3 题 4 分
%% typeId: 1
%% type: 单选题
%% chapter: 机械波
%% point: 波长频率波速
%% method: 无
%% score: 4
%% degree: 650
%% duplicateId: 0
%% body:
某仪器发射甲、乙两列横波，在同一均匀介质中相向传播，波速 $v$ 大小相等。某时刻的波形图如图所示，则这两列横波\xzanswer{C}

\onepicture[w=6.00cm]{figs/03.png}

\fourchoices[answer=C]
{在 $x = 9.0 \Um$ 处开始相遇}
{在 $x = 10 \Um$ 处开始相遇}
{波峰在 $x = 10.5 \Um$ 处相遇}
{波峰在 $x = 11.5 \Um$ 处相遇}

%% answer: C
%% memo:
\memoanswer{
AB．由题意可知两列波的波速相同，所以相同时间内传播的距离相同，故两列横波在 $x = 11.0 \Um$ 处开始相遇，故 AB 错误；

CD．甲波峰的坐标为 $x_{1} = 5 \Um$，乙波峰的坐标为 $x_{2} = 16 \Um$，由于两列波的波速相同，所以波峰在 $x' = 5 \Um + \frac{16 - 5}{2} \Um = 10.5 \Um$ 处相遇，故 C 正确，D 错误。

故选 C。
}

\item
%% number: 4
%% paperName: 2024年安徽高考第 4 题 4 分
%% typeId: 1
%% type: 单选题
%% chapter: 运动和力
%% point: 传送带
%% method: 无
%% score: 4
%% degree: 650
%% duplicateId: 0
%% body:
倾角为 $\theta$ 的传送带以恒定速率 $v_{0}$ 顺时针转动。$t = 0$ 时在传送带底端无初速轻放一小物块，如图所示。$t_{0}$ 时刻物块运动到传送带中间某位置，速度达到 $v_{0}$。不计空气阻力，则物块从传送带底端运动到顶端的过程中，加速度 $a$、速度 $v$ 随时间 $t$ 变化的关系图线可能正确的是\xzanswer{C}

\onepicture[w=5.00cm]{figs/04a.png}

\fourchoices[answer=C, ispicture=true, h=2.40cm]
{figs/04b1.jpg}
{figs/04b2.jpg}
{figs/04b3.jpg}
{figs/04b4.jpg}

%% answer: C
%% memo:
\memoanswer{
$0 \sim t_{0}$ 时间内：物体轻放在传送带上，做加速运动。受力分析可知，物体受重力、支持力、滑动摩擦力，滑动摩擦力大于重力的下滑分力，合力不变，故做匀加速运动。

$t_{0}$ 之后：当物块速度与传送带相同时，静摩擦力与重力的下滑分力相等，加速度突变为零，物块做匀速直线运动。

C 正确，ABD 错误。

故选 C。
}

\item
%% number: 5
%% paperName: 2024年安徽高考第 5 题 4 分
%% typeId: 1
%% type: 单选题
%% chapter: 曲线运动
%% point: 开普勒定律
%% method: 无
%% score: 4
%% degree: 650
%% duplicateId: 0
%% body:
2024 年 3 月 20 日，我国探月工程四期鹊桥二号中继星成功发射升空。当抵达距离月球表面某高度时，鹊桥二号开始进行近月制动，并顺利进入捕获轨道运行，如图所示，轨道的半长轴约为 $51\,900 \Ukm$。后经多次轨道调整，进入冻结轨道运行，轨道的半长轴约为 $9\,900 \Ukm$，周期约为 $24 \Uh$。则鹊桥二号在捕获轨道运行时\xzanswer{B}

\onepicture[w=5.00cm]{figs/05.png}

\fourchoices[answer=B]
{周期约为 $144 \Uh$}
{近月点的速度大于远月点的速度}
{近月点的速度小于在冻结轨道运行时近月点的速度}
{近月点的加速度大于在冻结轨道运行时近月点的加速度}

%% answer: B
%% memo:
\memoanswer{
A．冻结轨道和捕获轨道的中心天体是月球，根据开普勒第三定律 $\frac{T_{1}^{2}}{R_{1}^{3}} = \frac{T_{2}^{2}}{R_{2}^{3}}$

整理得 $T_{2} = \sqrt{\frac{R_{2}^{3}}{R_{1}^{3}}} T_{1} = 288 \Uh$

A 错误；

B．根据开普勒第二定律得，近月点的速度大于远月点的速度，B 正确；

C．近月点从捕获轨道到冻结轨道鹊桥二号进行近月制动，捕获轨道近月点的速度大于在冻结轨道运行时近月点的速度，C 错误；

D．两轨道的近月点所受的万有引力相同，根据牛顿第二定律可知，近月点的加速度等于在冻结轨道运行时近月点的加速度，D 错误。

故选 B。
}

\item
%% number: 6
%% paperName: 2024年安徽高考第 6 题 4 分
%% typeId: 1
%% type: 单选题
%% chapter: 运动和力
%% point: 变加速直线运动
%% method: 无
%% score: 4
%% degree: 550
%% duplicateId: 0
%% body:
如图所示，竖直平面内有两完全相同的轻质弹簧，它们的一端分别固定于水平线上的 M、N 两点，另一端均连接在质量为 $m$ 的小球上。开始时，在竖直向上的拉力作用下，小球静止于 MN 连线的中点 O，弹簧处于原长。后将小球竖直向上缓慢拉至 P 点，并保持静止，此时拉力 $F$ 大小为 $2mg$。已知重力加速度大小为 $g$，弹簧始终处于弹性限度内，不计空气阻力。若撤去拉力，则小球从 P 点运动到 O 点的过程中\xzanswer{A}

\onepicture[w=5.00cm]{figs/06.png}

\fourchoices[answer=A]
{速度一直增大}
{速度先增大后减小}
{加速度的最大值为 $3g$}
{加速度先增大后减小}

%% answer: A
%% memo:
\memoanswer{
AB．缓慢拉至 P 点，保持静止，由平衡条件可知此时拉力 $F$ 与重力和两弹簧的拉力合力为零。此时两弹簧的合力大小为 $mg$。当撤去拉力，则小球从 P 点运动到 O 点的过程中两弹簧的拉力与重力的合力始终向下，小球一直做加速运动，故 A 正确，B 错误；

CD．小球从 P 点运动到 O 点的过程中，形变量变小，弹簧在竖直方向的合力不断变小，故小球受的合外力一直变小，加速度的最大值为撤去拉力时的加速度，由牛顿第二定律可知

$2mg = ma$

加速度的最大值为 $2g$，CD 错误。

故选 A。
}

\item
%% number: 7
%% paperName: 2024年安徽高考第 7 题 4 分
%% typeId: 1
%% type: 单选题
%% chapter: 功和能
%% point: 功能原理
%% method: 建立模型
%% score: 4
%% degree: 550
%% duplicateId: 0
%% body:
在某地区的干旱季节，人们常用水泵从深水井中抽水灌溉农田，简化模型如图所示。水井中的水面距离水平地面的高度为 $H$。出水口距水平地面的高度为 $h$，与落地点的水平距离约为 $l$。假设抽水过程中 $H$ 保持不变，水泵输出能量的 $\eta$ 倍转化为水被抽到出水口处增加的机械能。已知水的密度为 $\rho$，水管内径的横截面积为 $S$，重力加速度大小为 $g$，不计空气阻力。则水泵的输出功率约为\xzanswer{B}

\onepicture[w=5.00cm]{figs/07.png}

\fourchoices[answer=B]
{$\frac{\rho gSl\sqrt{2gh}}{2\eta h}\left( H + h + \frac{l^{2}}{2h} \right)$}
{$\frac{\rho gSl\sqrt{2gh}}{2\eta h}\left( H + h + \frac{l^{2}}{4h} \right)$}
{$\frac{\rho gSl\sqrt{2gh}}{2\eta h}\left( H + \frac{l^{2}}{2h} \right)$}
{$\frac{\rho gSl\sqrt{2gh}}{2\eta h}\left( H + \frac{l^{2}}{4h} \right)$}

%% answer: B
%% memo:
\memoanswer{
设水从出水口射出的初速度为 $v_{0}$，取 $t$ 时间内的水为研究对象，该部分水的质量为 $m = v_{0}tS\rho$

根据平抛运动规律 $v_{0}t' = l$，$h = \frac{1}{2}gt'^{2}$

解得 $v_{0} = l\sqrt{\frac{g}{2h}}$

根据功能关系得

$Pt\eta = \frac{1}{2}mv_{0}^{2} + mg(H + h)$

联立解得水泵的输出功率为 $P = \frac{\rho gSl\sqrt{2gh}}{2\eta h}\left( H + h + \frac{l^{2}}{4h} \right)$

故选 B。
}

\item
%% number: 8
%% paperName: 2024年安徽高考第 8 题 4 分
%% typeId: 1
%% type: 单选题
%% chapter: 运动和力
%% point: 动量和能量
%% method: 无
%% score: 4
%% degree: 450
%% duplicateId: 0
%% body:
在某装置中的光滑绝缘水平面上，三个完全相同的带电小球，通过不可伸长的绝缘轻质细线，连接成边长为 $d$ 的正三角形，如图~\ref{2024安徽08a}所示。小球质量为 $m$，带电量为 $+q$，可视为点电荷。初始时，小球均静止，细线拉直。现将球 1 和球 2 间的细线剪断，当三个小球运动到同一条直线上时，速度大小分别为 $v_{1}$、$v_{2}$、$v_{3}$，如图~\ref{2024安徽08b}所示。该过程中三个小球组成的系统电势能减少了 $\frac{kq^{2}}{2d}$，$k$ 为静电力常量，不计空气阻力。则\xzanswer{D}

\twopicture[labelA=2024安徽08a, labelB=2024安徽08b, numA=甲, numB=乙, wA=3.30cm, wB=2.70cm]
{figs/08a.jpg}
{figs/08b.jpg}

\fourchoices[answer=D]
{该过程中小球 3 受到的合力大小始终不变}
{该过程中系统能量守恒，动量不守恒}
{在图乙位置，$v_{1} = v_{2}$，$v_{3} \neq 2v_{1}$}
{在图乙位置，$v_{3} = \sqrt{\frac{2kq^{2}}{3md}}$}

%% answer: D
%% memo:
\memoanswer{
AB．该过程中系统动能和电势能相互转化，能量守恒，对整个系统分析可知系统受到的合外力为 0，故动量守恒；当三个小球运动到同一条直线上时，根据对称性可知细线中的拉力相等，此时球 3 受到 1 和 2 的电场力大小相等，方向相反，故可知此时球 3 受到的合力为 0，球 3 从静止状态开始运动，瞬间受到的合力不为 0，故该过程中小球 3 受到的合力在改变，故 AB 错误；

CD．对系统根据动量守恒 $mv_{1} + mv_{2} = mv_{3}$

根据球 1 和 2 运动的对称性可知 $v_{1} = v_{2}$，解得 $v_{3} = 2v_{1}$

根据能量守恒

$\frac{1}{2}mv_{1}^{2} + \frac{1}{2}mv_{2}^{2} + \frac{1}{2}mv_{3}^{2} = \frac{kq^{2}}{2d}$

解得 $v_{3} = \sqrt{\frac{2kq^{2}}{3md}}$

故 C 错误，D 正确。

故选 D。
}

\item
%% number: 9
%% paperName: 2024年安徽高考第 9 题 5 分
%% typeId: 2
%% type: 多选题
%% chapter: 曲线运动
%% point: 类平抛运动
%% method: 无
%% score: 5
%% degree: 550
%% duplicateId: 0
%% body:
一倾角为 $30^{\circ}$ 足够大的光滑斜面固定于水平地面上，在斜面上建立 $Oxy$ 直角坐标系，如图~\ref{2024安徽09a}所示。从 $t = 0$ 开始，将一可视为质点的物块从 $O$ 点由静止释放，同时对物块施加沿 $x$ 轴正方向的力 $F_{1}$ 和 $F_{2}$，其大小与时间 $t$ 的关系如图~\ref{2024安徽09b}所示。已知物块的质量为 $1.2 \Ukg$，重力加速度 $g$ 取 $10 \Umsq$，不计空气阻力。则\xzanswer{BD}

\twopicture[labelA=2024安徽09a, labelB=2024安徽09b, numA=（1）, numB=（2）, wA=5.00cm, wB=3.60cm]
{figs/09a.jpg}
{figs/09b.jpg}

\fourchoices[answer=BD]
{物块始终做匀变速曲线运动}
{$t = 1 \Us$ 时，物块的 $y$ 坐标值为 $2.5 \Um$}
{$t = 1 \Us$ 时，物块的加速度大小为 $5\sqrt{3} \Umsq$}
{$t = 2 \Us$ 时，物块的速度大小为 $10\sqrt{2} \Ums$}

%% answer: BD
%% memo:
\memoanswer{
A．根据图像可得 $F_{1} = 4 - t$，$F_{2} = 3t$，故两力的合力为 $F = 4 + 2t$（单位：$\UN$）

物块在 $y$ 轴方向受到的力不变为 $mg\sin30^{\circ}$，$x$ 轴方向的力在改变，合力在改变，故物块做的不是匀变速曲线运动，故 A 错误；

B．在 $y$ 轴方向的加速度为

$a_{y} = \frac{mg\sin30^{\circ}}{m} = g\sin30^{\circ} = 5 \Umsq$

故 $t = 1 \Us$ 时，物块的 $y$ 坐标值为

$y = \frac{1}{2}a_{y}t^{2} = 2.5 \Um$

故 B 正确；

C．$t = 1 \Us$ 时，$F = 6 \UN$，故此时加速度大小为

$a = \sqrt{a_{x}^{2} + a_{y}^{2}} = \sqrt{\left(\frac{6}{1.2}\right)^{2} + 5^{2}} \Umsq = 5\sqrt{2} \Umsq$

故 C 错误；

D．对 $x$ 轴正方向，对物块根据动量定理

$Ft = mv_{x} - 0$

由于 $F$ 与时间 $t$ 成线性关系故可得

$\frac{(4 + 2 \times 0) + (4 + 2 \times 2)}{2} \times 2 = 1.2v_{x}$

解得 $v_{x} = 10 \Ums$

此时 $y$ 轴方向速度为

$v_{y} = g\sin30^{\circ} \cdot t = 5 \times 2 \Ums = 10 \Ums$

故此时物块的速度大小为 $v = \sqrt{v_{x}^{2} + v_{y}^{2}} = 10\sqrt{2} \Ums$

故 D 正确。

故选 BD。
}

\item
%% number: 10
%% paperName: 2024年安徽高考第 10 题 5 分
%% typeId: 2
%% type: 多选题
%% chapter: 磁场
%% point: 带电粒子在复合场中的运动
%% method: 无
%% score: 5
%% degree: 450
%% duplicateId: 0
%% body:
空间中存在竖直向下的匀强电场和垂直于纸面向里的匀强磁场，电场强度大小为 $E$，磁感应强度大小为 $B$。一质量为 $m$ 的带电油滴 a，在纸面内做半径为 $R$ 的圆周运动，轨迹如图所示。当 a 运动到最低点 P 时，瞬间分成两个小油滴 Ⅰ、Ⅱ，二者带电量、质量均相同。Ⅰ 在 P 点时与 a 的速度方向相同，并做半径为 $3R$ 的圆周运动，轨迹如图所示。Ⅱ 的轨迹未画出。已知重力加速度大小为 $g$，不计空气浮力与阻力以及 Ⅰ、Ⅱ 分开后的相互作用，则\xzanswer{ABD}

\onepicture[w=5.00cm]{figs/10.png}

\fourchoices[answer=ABD]
{油滴 a 带负电，所带电量的大小为 $\frac{mg}{E}$}
{油滴 a 做圆周运动的速度大小为 $\frac{gBR}{E}$}
{小油滴 Ⅰ 做圆周运动的速度大小为 $\frac{3gBR}{E}$，周期为 $\frac{4\pi E}{gB}$}
{小油滴 Ⅱ 沿顺时针方向做圆周运动}

%% answer: ABD
%% memo:
\memoanswer{
A．油滴 a 做圆周运动，故重力与电场力平衡，可知带负电，有 $mg = qE$

解得 $q = \frac{mg}{E}$

故 A 正确；

B．根据洛伦兹力提供向心力

$Bqv = m\frac{v^{2}}{R}$

得 $R = \frac{mv}{Bq}$

解得油滴 a 做圆周运动的速度大小为 $v = \frac{gBR}{E}$

故 B 正确；

C．设小油滴 Ⅰ 的速度大小为 $v_{1}$，得

$3R = \frac{\frac{m}{2}v_{1}}{B\frac{q}{2}}$

解得 $v_{1} = \frac{3BqR}{m} = \frac{3gBR}{E}$

周期为 $T = \frac{2\pi \cdot 3R}{v_{1}} = \frac{2\pi E}{gB}$

故 C 错误；

D．带电油滴 a 分离前后动量守恒，设分离后小油滴 Ⅱ 的速度为 $v_{2}$，取油滴 a 分离前瞬间的速度方向为正方向，得

$mv = \frac{m}{2}v_{1} + \frac{m}{2}v_{2}$

解得 $v_{2} = -\frac{gBR}{E}$

由于分离后的小油滴受到的电场力和重力仍然平衡，分离后小油滴 Ⅱ 的速度方向与正方向相反，根据左手定则可知小油滴 Ⅱ 沿顺时针方向做圆周运动，故 D 正确。

故选 ABD。
}

\item
%% number: 11
%% paperName: 2024年安徽高考第 11 题 0 分
%% typeId: 4
%% type: 实验题
%% chapter: 光学
%% point: 测定玻璃折射率
%% method: 无
%% score: 0
%% degree: 650
%% duplicateId: 0
%% body:
某实验小组做“测量玻璃的折射率”及拓展探究实验。
\begin{enumerate}
	\item
	为测量玻璃的折射率，按如图所示进行实验，以下表述正确的一项是\tkanswer{B}。（填正确答案标号）

	A．用笔在白纸上沿着玻璃砖上边和下边分别画出直线 a 和 $a'$

	B．在玻璃砖一侧插上大头针 $P_{1}$、$P_{2}$，眼睛在另一侧透过玻璃砖看两个大头针，使 $P_{2}$ 把 $P_{1}$ 挡住，这样就可以确定入射光线和入射点 $O_{1}$。在眼睛这一侧，插上大头针 $P_{3}$，使它把 $P_{1}$、$P_{2}$ 都挡住，再插上大头针 $P_{4}$，使它把 $P_{1}$、$P_{2}$、$P_{3}$ 都挡住，这样就可以确定出射光线和出射点 $O_{2}$

	C．实验时入射角 $\theta_{1}$ 应尽量小一些，以减小实验误差
	\item
	为探究介质折射率与光的频率的关系，分别用一束红光和一束绿光从同一点入射到空气与玻璃的分界面．保持相同的入射角，根据实验结果作出光路图，并标记红光和绿光，如图~\ref{2024安徽11b}所示．此实验初步表明：对于同一种介质，折射率与光的频率有关。频率大，折射率\tkanswer{大}（填“大”或“小”）
	\item
	为探究折射率与介质材料的关系，用同一束微光分别入射玻璃砖和某透明介质，如图~\ref{2024安徽11c}、\ref{2024安徽11d}所示。保持相同的入射角 $\alpha_{1}$，测得折射角分别为 $\alpha_{2}$、$\alpha_{3}$（$\alpha_{2} < \alpha_{3}$），则玻璃和该介质的折射率大小关系为 $n_{\text{玻璃}}$\tkanswer{$>$}$n_{\text{介质}}$（填“$>$”或“$<$”）。此实验初步表明：对于一定频率的光，折射率与介质材料有关。
\end{enumerate}
\onepicture[w=3.20cm, align=right]{figs/11a.png}
\threepicture[labelA=2024安徽11b, labelB=2024安徽11c, labelC=2024安徽11d, numA=乙, numB=丙, numC=丁, wA=3.00cm, wB=3.00cm, wC=3.00cm, align=right]
{figs/11b.jpg}
{figs/11c.jpg}
{figs/11d.jpg}

%% answer:
\jdanswer{
\begin{enumerate}
	\item
	B

	\item
	大

	\item
	$>$

\end{enumerate}
}
%% memo:
\memoanswer{
（1）A．在白纸上画出一条直线 a 作为界面，把长方体玻璃砖放在白纸上，使它的一个长边与 a 对齐。用直尺或者三角板轻靠在玻璃砖的另一长边，按住直尺或三角板不动，将玻璃砖取下，画出直线 $a'$ 代表玻璃砖的另一边，而不能用笔在白纸上沿着玻璃砖上边和下边分别画出直线 a 和 $a'$，故 A 错误；

B．在玻璃砖一侧插上大头针 $P_{1}$、$P_{2}$，眼睛在另一侧透过玻璃砖看两个大头针，使 $P_{2}$ 把 $P_{1}$ 挡住，这样就可以确定入射光线和入射点 $O_{1}$。在眼睛这一侧，插上大头针 $P_{3}$，使它把 $P_{1}$、$P_{2}$ 都挡住，再插上大头针 $P_{4}$，使它把 $P_{1}$、$P_{2}$、$P_{3}$ 都挡住，这样就可以确定出射光线和出射点 $O_{2}$，故 B 正确；

C．实验时入射角 $\theta_{1}$ 应尽量大一些，但也不能太大（接近 $90^{\circ}$），以减小实验误差，故 C 错误。

故选 B。

（2）由图乙可知，入射角相同，绿光的折射角小于红光的折射角，根据光的折射定律

$n = \frac{\sin\alpha}{\sin\beta}$

可知绿光的折射率大于红光的折射率，又因为绿光的频率大于红光的频率，所以频率大，折射率大。

（3）根据折射定律可知，玻璃的折射率为

$n_{\text{玻璃}} = \frac{\sin\alpha_{1}}{\sin\alpha_{2}}$

该介质的折射率为

$n_{\text{介质}} = \frac{\sin\alpha_{1}}{\sin\alpha_{3}}$

其中 $\alpha_{2} < \alpha_{3}$，所以 $n_{\text{玻璃}} > n_{\text{介质}}$。
}

\item
%% number: 12
%% paperName: 2024年安徽高考第 12 题 10 分
%% typeId: 1
%% type: 单选题
%% chapter: 电路
%% point: 电路实验
%% method: 无
%% score: 10
%% degree: 950
%% duplicateId: 0
%% body:
某实验小组要将电流表 G（铭牌标示：$I_{g} = 500 \UuA$，$R_{g} = 800 \UO$）改装成量程为 $1 \UV$ 和 $3 \UV$ 的电压表，并用标准电压表对其进行校准。选用合适的电源、滑动变阻器、电阻箱、开关和标准电压表等实验器材，按图~\ref{2024安徽12a}所示连接电路，其中虚线框内为改装电路。
\begin{enumerate}
	\item
	开关 $S_{1}$ 闭合前，滑片 P 应移动到\tkanswer{M}（填“M”或“N”）端。
	\item
	根据要求和已知信息，电阻箱 $R_{1}$ 的阻值已调至 $1200 \UO$，则 $R_{2}$ 的阻值应调至\tkanswer{4000}$\UO$。
	\item
	当单刀双掷开关 $S_{2}$ 与 a 连接时，电流表 G 和标准电压表 V 的示数分别为 $I$、$U$，则电流表 G 的内阻可表示为\tkanswer{$R_{g} = \frac{U}{I} - R_{1} - R_{2}$}。（结果用 $U$、$I$、$R_{1}$、$R_{2}$ 表示）
	\item
	校准电表时，发现改装后电压表的读数始终比标准电压表的读数偏大，经排查发现电流表 G 内阻的真实值与铭牌标示值有偏差，则只要\tkanswer{A}即可。（填正确答案标号）

	A．增大电阻箱 $R_{1}$ 的阻值

	B．减小电阻箱 $R_{2}$ 的阻值

	C．将滑动变阻器的滑片 P 向 M 端滑动
	\item
	校准完成后，开关 $S_{2}$ 与 b 连接，电流表 G 的示数如图~\ref{2024安徽12b}所示，此示数对应的改装电压表读数为\tkanswer{0.86}$\UV$。（保留 2 位有效数字）
\end{enumerate}
\twopicture[labelA=2024安徽12a, labelB=2024安徽12b, numA=（1）, numB=（2）, wA=5.00cm, wB=5.00cm]
{figs/12a.jpg}
{figs/12b.jpg}

%% answer: IA
\jdanswer{
\begin{enumerate}
	\item
	M

	\item
	4000

	\item
	$R_{g} = \frac{U}{I} - R_{1} - R_{2}$

	\item
	A

	\item
	0.86

\end{enumerate}
}
%% memo:
\memoanswer{
（1）由图可知，该滑动变阻器采用分压式接法，为了保护电路，在开关 $S_{1}$ 闭合前，滑片 P 应移到 M 端；

（2）当开关 $S_{2}$ 接 b 时，电压表量程为 $1 \UV$，根据欧姆定律

$U_{1} = I_{g}(R_{g} + R_{1})$

当开关 $S_{2}$ 接 a 时，电压表量程为 $3 \UV$，根据欧姆定律

$U_{2} = I_{g}(R_{g} + R_{1} + R_{2})$

其中 $R_{1} = 1200 \UO$

联立解得 $R_{2} = 4000 \UO$

（3）当开关 $S_{2}$ 接 a 时，根据欧姆定律 $U = I(R_{g} + R_{1} + R_{2})$

可得电流表 G 的内阻可表示为 $R_{g} = \frac{U}{I} - R_{1} - R_{2}$

（4）校准电表时，发现改装后电压表的读数始终比标准电压表的读数偏大，可知电流表 G 内阻的真实值小于铭牌标示值，根据闭合电路的欧姆定律可以增大两电阻箱的阻值。

故选 A。

（5）根据闭合电路欧姆定律 $U_{V} = I_{A}(R_{g} + R_{1}) = 430 \times 10^{-6} \times (800 + 1200) \UV = 0.86 \UV$。
}

\item
%% number: 13
%% paperName: 2024年安徽高考第 13 题 0 分
%% typeId: 6
%% type: 计算题
%% chapter: 气体性质
%% point: 玻意耳定律
%% method: 无
%% score: 0
%% degree: 750
%% duplicateId: 0
%% body:
某人驾驶汽车，从北京到哈尔滨，在哈尔滨发现汽车的某个轮胎内气体的压强有所下降（假设轮胎内气体的体积不变，且没有漏气，可视为理想气体）。于是在哈尔滨给该轮胎充入压强与大气压相同的空气，使其内部气体的压强恢复到出发时的压强（假设充气过程中，轮胎内气体的温度与环境相同，且保持不变）。已知该轮胎内气体的体积 $V_{0} = 30 \UL$，从北京出发时，该轮胎气体的温度 $t_{1} = -3^{\circ}C$，压强 $p_{1} = 2.7 \times 10^{5} \UPa$。哈尔滨的环境温度 $t_{2} = -23^{\circ}C$，大气压强 $p_{0}$ 取 $1.0 \times 10^{5} \UPa$。求：
\begin{enumerate}
	\item
	在哈尔滨时，充气前该轮胎气体压强的大小。
	\item
	充进该轮胎的空气体积。
\end{enumerate}

%% answer:
\jdanswer{
\begin{enumerate}
	\item
	$2.5 \times 10^{5} \UPa$

	\item
	$6 \UL$

\end{enumerate}
}
%% memo:
\memoanswer{
（1）由查理定律可得

$\frac{p_{1}}{T_{1}} = \frac{p_{2}}{T_{2}}$

其中 $p_{1} = 2.7 \times 10^{5} \UPa$，$T_{1} = 273 - 3\ \UK = 270 \UK$，$T_{2} = 273 - 23\ \UK = 250 \UK$

代入数据解得，在哈尔滨时，充气前该轮胎气体压强的大小为

$p_{2} = 2.5 \times 10^{5} \UPa$

（2）由玻意耳定律 $p_{2}V_{0} + p_{0}V = p_{1}V_{0}$

代入数据解得，充进该轮胎的空气体积为 $V = 6 \UL$。
}

\item
%% number: 14
%% paperName: 2024年安徽高考第 14 题 14 分
%% typeId: 6
%% type: 计算题
%% chapter: 运动和力
%% point: 动量和能量
%% method: 无
%% score: 14
%% degree: 450
%% duplicateId: 0
%% body:
如图所示，一实验小车静止在光滑水平面上，其上表面有粗糙水平轨道与光滑四分之一圆弧轨道。圆弧轨道与水平轨道相切于圆弧轨道最低点，一物块静止于小车最左端，一小球用不可伸长的轻质细线悬挂于 O 点正下方，并轻靠在物块右侧。现将细线拉直到水平位置时，静止释放小球，小球运动到最低点时与物块发生弹性碰撞。碰撞后，物块沿着的轨道运动，已知细线长 $L = 1.25 \Um$。小球质量 $m = 0.20 \Ukg$。物块、小车质量均为 $M = 0.30 \Ukg$。小车上的水平轨道长 $s = 1.0 \Um$。圆弧轨道半径 $R = 0.15 \Um$。小球、物块均可视为质点。不计空气阻力，重力加速度 $g$ 取 $10 \Umsq$。
\begin{enumerate}
	\item
	求小球运动到最低点与物块碰撞前所受拉力的大小；
	\item
	求小球与物块碰撞后的瞬间，物块速度的大小；
	\item
	为使物块能进入圆弧轨道，且在上升阶段不脱离小车，求物块与水平轨道间的动摩擦因数 $\mu$ 的取值范围。
\end{enumerate}
\onepicture[w=5.50cm, align=right]{figs/14.png}

%% answer:
\jdanswer{
\begin{enumerate}
	\item
	$6 \UN$

	\item
	$4 \Ums$

	\item
	$0.25 \leqslant \mu < 0.4$

\end{enumerate}
}
%% memo:
\memoanswer{
（1）对小球摆动到最低点的过程中，由动能定理

$mgL = \frac{1}{2}mv_{0}^{2} - 0$

解得 $v_{0} = 5 \Ums$

在最低点，对小球由牛顿第二定律

$F_{T} - mg = m\frac{v_{0}^{2}}{L}$

解得，小球运动到最低点与物块碰撞前所受拉力的大小为 $F_{T} = 6 \UN$

（2）小球与物块碰撞过程中，由动量守恒定律和机械能守恒定律

$mv_{0} = mv_{1} + Mv_{2}$

$\frac{1}{2}mv_{0}^{2} = \frac{1}{2}mv_{1}^{2} + \frac{1}{2}Mv_{2}^{2}$

解得小球与物块碰撞后的瞬间，物块速度的大小为

$v_{2} = \frac{2m}{m + M}v_{0} = 4 \Ums$

（3）若物块恰好运动到圆弧轨道的最低点，此时两者共速，则对物块与小车整体由水平方向动量守恒

$Mv_{2} = 2Mv_{3}$

由能量守恒定律

$\frac{1}{2}Mv_{2}^{2} = \frac{1}{2} \times 2Mv_{3}^{2} + \mu_{1}Mgs$

解得 $\mu_{1} = 0.4$

若物块恰好运动到与圆弧圆心等高的位置，此时两者共速，则对物块与小车整体由水平方向动量守恒

$Mv_{2} = 2Mv_{4}$

由能量守恒定律

$\frac{1}{2}Mv_{2}^{2} = \frac{1}{2} \times 2Mv_{4}^{2} + \mu_{2}Mgs + MgR$

解得 $\mu_{2} = 0.25$

综上所述物块与水平轨道间的动摩擦因数 $\mu$ 的取值范围为 $0.25 \leqslant \mu < 0.4$。
}

\item
%% number: 15
%% paperName: 2024年安徽高考第 15 题 18 分
%% typeId: 6
%% type: 计算题
%% chapter: 电磁感应
%% point: 电磁感应中的匀变速
%% method: 极值
%% score: 18
%% degree: 350
%% duplicateId: 0
%% body:
如图所示，一“U”型金属导轨固定在竖直平面内，一电阻不计，质量为 $m$ 的金属棒 ab 垂直于导轨，并静置于绝缘固定支架上。边长为 $L$ 的正方形 cdef 区域内，存在垂直于纸面向外的匀强磁场。支架上方的导轨间，存在竖直向下的匀强磁场。两磁场的磁感应强度大小 $B$ 随时间的变化关系均为 $B = kt$（SI），$k$ 为常数（$k > 0$）。支架上方的导轨足够长，两边导轨单位长度的电阻均为 $r$，下方导轨的总电阻为 $R$。$t = 0$ 时，对 ab 施加竖直向上的拉力，恰使其向上做加速度大小为 $a$ 的匀加速直线运动，整个运动过程中 ab 与两边导轨接触良好。已知 ab 与导轨间动摩擦因数为 $\mu$，重力加速度大小为 $g$。不计空气阻力，两磁场互不影响。
\begin{enumerate}
	\item
	求通过面积 $S_{cdef}$ 的磁通量大小随时间 $t$ 变化的关系式，以及感应电动势的大小，并写出 ab 中电流的方向；
	\item
	求 ab 所受安培力的大小随时间 $t$ 变化的关系式；
	\item
	求经过多长时间，对 ab 所施加的拉力达到最大值，并求此最大值。
\end{enumerate}
\onepicture[w=4.50cm, align=right]{figs/15.png}

%% answer:
\jdanswer{
\begin{enumerate}
	\item
	$kL^{2}t$，$kL^{2}$，从 a 流向 b

	\item
	$F_{\text{安}} = \frac{k^{2}L^{3}t}{R + art^{2}}$

	\item
	$\frac{\mu k^{2}L^{3}}{2\sqrt{Rar}} + m(g + a)$

\end{enumerate}
}
%% memo:
\memoanswer{
（1）通过面积 $S_{cdef}$ 的磁通量大小随时间 $t$ 变化的关系式为

$\Phi = BS = kL^{2}t$

根据法拉第电磁感应定律得

$E = n\frac{\Delta\Phi}{\Delta t} = \frac{\Delta B \cdot S}{\Delta t} = kL^{2}$

由楞次定律可知 ab 中的电流从 a 流向 b。

（2）根据左手定则可知 ab 受到的安培力方向垂直导轨面向里，大小为 $F_{\text{安}} = BIL$

其中 $B = kt$

设金属棒向上运动的位移为 $x$，则根据运动学公式 $x = \frac{1}{2}at^{2}$

所以导轨上方的电阻为 $R' = 2xr$

由闭合电路欧姆定律得 $I = \frac{E}{R + 2xr}$

联立得 ab 所受安培力的大小随时间 $t$ 变化的关系式为

$F_{\text{安}} = \frac{k^{2}L^{3}t}{R + art^{2}}$

（3）由题知 $t = 0$ 时，对 ab 施加竖直向上的拉力，恰使其向上做加速度大小为 $a$ 的匀加速直线运动，则对 ab 受力分析由牛顿第二定律

$F - mg - \mu F_{\text{安}} = ma$

其中 $F_{\text{安}} = \frac{k^{2}L^{3}t}{R + art^{2}}$

联立可得

$F = \frac{\mu k^{2}L^{3}t}{R + art^{2}} + m(g + a) = \frac{\mu k^{2}L^{3}}{\frac{R}{t} + art} + m(g + a)$

根据均值不等式可知，当 $\frac{R}{t} = art$ 时，$F$ 有最大值，故解得 $t = \sqrt{\frac{R}{ar}}$

$F$ 的最大值为 $F_{m} = \frac{\mu k^{2}L^{3}}{2\sqrt{Rar}} + m(g + a)$。
}

\end{enumerate}

\end{document}
